\chapter{PS/2 Keyboard}
\label{cap:keyboard}

\section{Overview}

The \code{src/keyboard.rs} driver implements the subset needed to read keys
from a PS/2 keyboard connected to the 8042 controller. It combines concepts
from two Linux subsystems:

\begin{itemize}
  \item \code{drivers/input/serio/i8042.c}: controller access and the command
        protocol;
  \item \code{drivers/input/keyboard/atkbd.c}: scancode and modifier key
        interpretation.
\end{itemize}

There is no IRQ1, event queue, \code{input\_device}, Linux \emph{keycode}, or
\code{evdev}. The driver does \emph{polling} of the 8042 output buffer and
returns either an ASCII character or a console switch request to the TTY
layer.

\section{Hardware: 8042 controller}

The 8042 exposes two ports:

\begin{table}[H]
\centering
\caption{Ports of the 8042 controller.}
\label{tab:ports8042}
\begin{tabular}{@{}lll@{}}
\toprule
\textbf{Port} & \textbf{Name} & \textbf{Usage} \\
\midrule
\code{0x60} & Data & Reading scancodes; writing data to the keyboard. \\
\code{0x64} & Status/Command & Reading status; writing commands to the 8042. \\
\bottomrule
\end{tabular}
\end{table}

The status register (\code{0x64}) has the bits used by the driver:

\begin{table}[H]
\centering
\caption{8042 status bits.}
\label{tab:status8042}
\begin{tabular}{@{}lll@{}}
\toprule
\textbf{Bit} & \textbf{Mask} & \textbf{Meaning} \\
\midrule
0 & \code{0x01} & Output buffer full (there is data to read) \\
1 & \code{0x02} & Input buffer full (do not send a command) \\
5 & \code{0x20} & Data came from the auxiliary device (mouse) \\
\bottomrule
\end{tabular}
\end{table}

\section{Initialization}

Initialization is more than clearing the buffer. The BIOS or GRUB may leave
the controller in a state where the keyboard sends set 2 scancodes, while the
driver expects set 1, or with scanning disabled. The procedure implemented is:

\begin{table}[H]
\centering
\caption{8042 initialization sequence.}
\label{tab:kbdinit}
\begin{tabular}{@{}clp{7.3cm}@{}}
\toprule
\textbf{\#} & \textbf{Command} & \textbf{Effect} \\
\midrule
1 & reads from \code{0x60} & Discards pending bytes in the output buffer. \\
2 & \code{0xAD} in \code{0x64} & Disables the keyboard port. \\
3 & \code{0x20} in \code{0x64} & Requests a read of the \emph{configuration byte}. \\
4 & read from \code{0x60} & Obtains the current configuration byte
(or uses \code{0x45} if there is no response). \\
5 & adjustment & Turns on translation (bit 6), turns off IRQ1 (bit 0), and
turns on the clock (clears bit 4). \\
6 & \code{0x60} in \code{0x64} & Prepares the write of the new configuration
byte. \\
7 & byte in \code{0x60} & Sends the new \emph{configuration byte}. \\
8 & \code{0xAE} in \code{0x64} & Re-enables the keyboard port. \\
9 & \code{0xF4} in \code{0x60} & Sends the \emph{enable scanning} command to
the keyboard. \\
10 & read from \code{0x60} & Consumes the ACK (\code{0xFA}), if present. \\
11 & cleanup & Discards any residue from the buffer. \\
\bottomrule
\end{tabular}
\end{table}

\begin{nota}[Why translate?]
The modern PS/2 keyboard usually speaks \emph{scan code set 2}, but the 8042
controller can present \emph{set 1} to the system when the translation bit is
on. Linux also configures this translation in i8042. Without it, keys arrive in
codes different from those expected by the driver's table.
\end{nota}

The waits use up to 100,000 iterations, both for the input buffer to empty and
for the output buffer to fill. If the controller does not respond,
initialization continues instead of hanging.

\section{Scancodes and ASCII translation}

The driver works with set 1. The \code{KEYMAP} table has 128 positions and
covers letters, digits, space, and some symbols. Unmapped positions are zero
and are ignored.

\begin{table}[H]
\centering
\caption{Examples of set 1 scancodes.}
\label{tab:scancodes}
\begin{tabular}{@{}llll@{}}
\toprule
\textbf{Key} & \textbf{Make} & \textbf{Break} & \textbf{Result} \\
\midrule
\code{A} & \code{0x1E} & \code{0x9E} & \code{a} / \code{A} with Shift \\
\code{1} & \code{0x02} & \code{0x82} & \code{1} / \code{!} with Shift \\
\code{Enter} & \code{0x1C} & \code{0x9C} & \code{\textbackslash n} \\
\code{Backspace} & \code{0x0E} & \code{0x8E} & \code{0x08} \\
\code{Tab} & \code{0x0F} & \code{0x8F} & \code{\textbackslash t} \\
\code{Space} & \code{0x39} & \code{0xB9} & space \\
\code{Left Shift} & \code{0x2A} & \code{0xAA} & enables/disables Shift \\
\code{Right Shift} & \code{0x36} & \code{0xB6} & enables/disables Shift \\
\code{Left Alt} & \code{0x38} & \code{0xB8} & enables/disables Alt \\
\code{F1}..\code{F8} & \code{0x3B}..\code{0x42} & \code{0xBB}..\code{0xC2}
& \code{Alt+F1}..\code{Alt+F8} selects \code{tty0}..\code{tty7} \\
\code{Extended key} & \code{0xE0} & --- & prefix; next code ignored \\
\bottomrule
\end{tabular}
\end{table}

\subsection{Shift}

The Shift state is kept in \code{static SHIFT: AtomicBool}. The \code{0x2A}
and \code{0x36} keys turn the state on; the \emph{break codes} \code{0xAA} and
\code{0xB6} turn it off. The \code{apply\_shift} function transforms letters to
uppercase and applies the secondary symbol of digits and punctuation:

\begin{lstlisting}[caption={Mapping with Shift.},label={lst:shift}]
b'a'..=b'z' => byte - 32,
b'1'..=b'9' => b"!@#$%^&*("[(byte - b'1') as usize],
b'0' => b')',
b'-' => b'_',
b'=' => b'+',
/* ... */
\end{lstlisting}

\subsection{Extended keys}

The \code{0xE0} prefix sets the \code{EXTENDED} flag. The next scancode is
discarded and the flag is cleared. This prevents arrows, Insert, Delete, Home,
End, and numeric keypad keys from being interpreted as letters. The
consequence is that these keys have no function in the shell.

\subsection{Alt and console selection}

The Alt state is kept in \code{static ALT: AtomicBool}: scancode \code{0x38}
sets it and the break code \code{0xB8} clears it. While Alt is held, the F1
through F8 scancodes (\code{0x3B}..\code{0x42}) are not translated to ASCII.
They produce \code{Key::Console(0..7)}, which the TTY layer uses to switch to
\code{tty0}..\code{tty7}.

\begin{lstlisting}[caption={Result of the keyboard driver.},label={lst:key}]
pub enum Key {
    Char(u8),
    Console(usize),
}
\end{lstlisting}

\subsection{Mouse}

Mouse data shares the same 8042. If bit 5 of the status is on, the byte read
is discarded. This prevents mouse movements from appearing as letters.

\section{Read flow}

\begin{figure}[H]
\centering
\begin{tikzpicture}[
  kstep/.style={draw=rinuxblue, fill=rinuxlight, rounded corners=2pt,
               minimum width=5.1cm, minimum height=0.8cm,
               align=center, font=\scriptsize},
  dec/.style={draw=orange!70!black, fill=orange!8, diamond, aspect=2.2,
              inner sep=1pt, align=center, font=\scriptsize},
  arrow/.style={-{Latex[length=2mm]}, thick, rinuxblue},
  node distance=0.38cm
]
\node[kstep] (status) {Reads status (\code{0x64})};
\node[dec, below=of status] (full) {bit 0 set?};
\node[kstep, below=of full] (read) {Reads data (\code{0x60})};
\node[dec, below=of read] (mouse) {bit 5 (mouse)?};
\node[dec, below=of mouse] (ext) {\code{0xE0} or extended?};
\node[dec, below=of ext] (break) {bit 7 (break)?};
\node[dec, below=of break] (alt) {Alt+F1..F8?};
\node[kstep, below=of alt] (key) {Translates and returns \code{Key}};
\node[kstep, right=2.4cm of full] (none) {Returns \code{None}};
\draw[arrow] (status) -- (full);
\draw[arrow] (full) -- node[right,font=\tiny]{yes} (read);
\draw[arrow] (full.east) -- node[above,font=\tiny]{no} (none.west);
\draw[arrow] (read) -- (mouse);
\draw[arrow] (mouse) -- node[right,font=\tiny]{no} (ext);
\draw[arrow] (ext) -- node[right,font=\tiny]{no} (break);
\draw[arrow] (break) -- node[right,font=\tiny]{no} (alt);
\draw[arrow] (alt) -- node[right,font=\tiny]{no} (key);
\draw[arrow] (alt.east) -- node[above,font=\tiny]{yes} ++(0.65,0)
  node[right,font=\tiny]{\code{Console(n)}};
\draw[arrow] (mouse.east) -- ++(0.9,0) |- (none.east);
\draw[arrow] (ext.east) -- ++(0.5,0) |- (none.east);
\draw[arrow] (break.east) -- ++(0.2,0) |- (none.east);
\end{tikzpicture}
\caption{Flow of \code{read\_key()}. Discard paths lead to \code{None}.}
\label{fig:kbdflow}
\end{figure}

The TTY layer calls \code{keyboard::read\_key()} repeatedly inside
\code{tty::pump\_input}. A \code{Key::Char} is pushed to the input queue of the
active console; a \code{Key::Console(n)} switches the active console. The
compatibility wrapper \code{read\_char()} discards console requests and
returns only characters.

\section{API and internal functions}

\begin{table}[H]
\centering
\caption{Functions of \code{keyboard.rs}.}
\label{tab:kbdapi}
\begin{tabularx}{\linewidth}{@{}lX@{}}
\toprule
\textbf{Function} & \textbf{Description} \\
\midrule
\code{init()} & Configures 8042 translation, IRQ, clock, and scanning. \\
\code{read\_key()} & Reads and translates a scancode; returns
\code{Option<Key>}. \\
\code{read\_char()} & Compatibility wrapper; returns only \code{Key::Char}. \\
\code{build\_keymap()} & Builds the 128-entry table at compile time. \\
\code{apply\_shift(byte)} & Applies the effect of Shift to a character. \\
\code{wait\_input\_clear()} & Waits for the 8042 input buffer to empty. \\
\code{wait\_output\_full()} & Waits for a byte to become available at the output. \\
\code{flush\_output()} & Discards up to 64 pending bytes. \\
\code{send\_command(cmd)} & Sends a command to the 8042 (port \code{0x64}). \\
\code{send\_data(data)} & Sends a byte to the keyboard (port \code{0x60}). \\
\bottomrule
\end{tabularx}
\end{table}

\section{Limitations}

\begin{itemize}
  \item \textbf{Polling}: there is no IRQ1; the shell must poll constantly. The
        CPU stays busy even without input.
  \item \textbf{No Caps Lock, Num Lock, or Scroll Lock}: the LEDs and states
        are not handled.
  \item \textbf{No AltGr}: layouts that depend on AltGr do not produce
        additional characters.
  \item \textbf{No functional extended keys}: arrows, Delete, Home, End,
        Page Up/Down, and the numeric keypad are discarded.
  \item \textbf{No ABNT2 layout}: the table is based on US/QWERTY; the
        physical key positions and the symbol set may differ.
  \item \textbf{No useful \emph{key release}}: apart from Shift and Alt, the
        \emph{break codes} are ignored; there is no key press detection for
        games or shortcuts.
  \item \textbf{No error handling}: error bytes (\code{0x00}, \code{0xFF},
        \code{0xEE}) have no specific handling.
  \item \textbf{Hardware dependency}: assumes an 8042-compatible device. On
        USB-HID-only systems, the keyboard may not work without \emph{legacy
        USB} enabled in the BIOS.
\end{itemize}
